\documentclass[10pt,a4paper]{amsart}
\usepackage{amsmath,amssymb,amsthm,graphicx}
% Retained preamble for the Stacks Project chapter snapshot in this directory.
% The chapter's own preamble.tex is a copy of the shared Stacks Project preamble; the block
% below is sliced MECHANICALLY, line for line, from it. Nothing was added, removed or
% redefined; the wrapper class, the xr-hyper/hyperref wiring and the externaldocument list
% are omitted because the assembler supplies the document class.
%
% lines taken from preamble.tex: 3, 5, 9, 11, 12, 156, 157, 158, ... (51 lines total)
%
% package coverage restored from the chapter's own preamble (2026-10-13): xy, xy-2cell

\usepackage[all]{xy}
\xyoption{2cell}
\UseAllTwocells

\newenvironment{reference}{\comment}{\endcomment}
\newenvironment{slogan}{\comment}{\endcomment}
\newenvironment{history}{\comment}{\endcomment}
\theoremstyle{plain}
\newtheorem{theorem}[subsection]{Theorem}
\newtheorem{proposition}[subsection]{Proposition}
\newtheorem{lemma}[subsection]{Lemma}
\theoremstyle{definition}
\newtheorem{definition}[subsection]{Definition}
\newtheorem{example}[subsection]{Example}
\newtheorem{exercise}[subsection]{Exercise}
\newtheorem{situation}[subsection]{Situation}
\theoremstyle{remark}
\newtheorem{remark}[subsection]{Remark}
\newtheorem{remarks}[subsection]{Remarks}
\numberwithin{equation}{subsection}
\def\lim{\mathop{\mathrm{lim}}\nolimits}
\def\colim{\mathop{\mathrm{colim}}\nolimits}
\def\Spec{\mathop{\mathrm{Spec}}}
\def\Hom{\mathop{\mathrm{Hom}}\nolimits}
\def\Ext{\mathop{\mathrm{Ext}}\nolimits}
\def\SheafHom{\mathop{\mathcal{H}\!\mathit{om}}\nolimits}
\def\SheafExt{\mathop{\mathcal{E}\!\mathit{xt}}\nolimits}
\def\Sch{\mathit{Sch}}
\def\Mor{\mathop{\mathrm{Mor}}\nolimits}
\def\Ob{\mathop{\mathrm{Ob}}\nolimits}
\def\Sh{\mathop{\mathit{Sh}}\nolimits}
\def\NL{\mathop{N\!L}\nolimits}
\def\CH{\mathop{\mathrm{CH}}\nolimits}
\def\proetale{{pro\text{-}\acute{e}tale}}
\def\etale{{\acute{e}tale}}
\def\QCoh{\mathit{QCoh}}
\def\Ker{\mathop{\mathrm{Ker}}}
\def\Im{\mathop{\mathrm{Im}}}
\def\Coker{\mathop{\mathrm{Coker}}}
\def\Coim{\mathop{\mathrm{Coim}}}
\DeclareMathSymbol{\boxtimes}{\mathbin}{AMSa}{"02}
\def\QCohstack{\mathcal{QC}\!\mathit{oh}}
\def\Cohstack{\mathcal{C}\!\mathit{oh}}
\def\Spacesstack{\mathcal{S}\!\mathit{paces}}
\def\Quotfunctor{\mathrm{Quot}}
\def\Hilbfunctor{\mathrm{Hilb}}
\def\Curvesstack{\mathcal{C}\!\mathit{urves}}
\def\Polarizedstack{\mathcal{P}\!\mathit{olarized}}
\def\Complexesstack{\mathcal{C}\!\mathit{omplexes}}
\def\Pic{\mathop{\mathrm{Pic}}\nolimits}
\def\Picardstack{\mathcal{P}\!\mathit{ic}}
\def\Picardfunctor{\mathrm{Pic}}
\def\Deformationcategory{\mathcal{D}\!\mathit{ef}}
\usepackage{cleveref}
\crefname{originallemma}{original lemma}{original lemmas}
\Crefname{originallemma}{Original lemma}{Original lemmas}
\makeatletter
\newlabel{lemma-difference-divisors-on-curves}{{0B9F}{source}{}{sourcepointer-1}{}}
\newlabel{lemma-difference-divisors-on-curves@cref}{{[originallemma][0][]0B9F}{[1][0][]source}{}{}{}}
\makeatother
\title{Adding a point is finite locally free of degree d+1}
\author{The Stacks Project Authors}
\date{stacks-project:pic.tex:local-snapshot-20261012; GNU FDL 1.2 or later}
\begin{document}
\maketitle
\noindent\textit{Editorial note.} Statement span is lines 785-792 of sources/stacks-pic/pinned-original-source.tex; proof span is lines 794-816, the complete author proof (21 body lines).\par\medskip
\noindent\textit{Editorial note.} Navigation stubs: 1 cited-but-not-included labels are supplied as non-load-bearing pointers with their pinned tags where the snapshot tag index has them.\par\medskip
\noindent\textit{Editorial note.} extra\_packages is empty on purpose: the retained preamble already loads xy with options and a second declaration clashes in the assembler wrapper.\par\medskip
\noindent\textit{Editorial note.} Text is retained verbatim, including the author's own typos and placeholders.\par\medskip
\noindent\textit{Editorial note.} Queue index 15 of queue\_pic.json (pb=986, stmt 785-792, proof 794-816).\par\medskip
\noindent\textit{Editorial note.} stacks-pic wave-1 (pi01): fresh adjudication accept (2026-10-14); pool per standing queue hypercovering -> pic (SMC reassigned to Zhen per Chief 08:5x); pic.tex 27 labels/23 with proof judged 21 accept / 2 reject / 4 pb=0 env-dup rows excluded; spans cite-free and reference/slogan/history-free so no patchenv needed.\par\medskip
\noindent\textit{Allowed standard prerequisites (editorial).}\begin{itemize}
\item Explicit prerequisite boundary. The cited inputs are prerequisites; the argument in the retained span is not replaced by any cited result.
\end{itemize}\medskip
\ifdefined\hypertarget\hypertarget{sourcepointer-1}{}\fi
\noindent\textit{Source reference.} Original lemma 0B9F: sources/stacks-pic/pinned-original-source.tex, source lines 726--732. Cited by the retained statement or proof but not part of the collected spans; supplied only as a navigation pointer. The argument does not depend on its content.\par\medskip
\begin{lemma}
\label{lemma-universal-object}
Let $X \to S$ be a smooth morphism of schemes of relative dimension $1$
such that the functors $\Hilbfunctor^d_{X/S}$ are representable. The morphism
$\underline{\Hilbfunctor}^d_{X/S} \times_S X \to
\underline{\Hilbfunctor}^{d + 1}_{X/S}$
is finite locally free of degree $d + 1$.
\end{lemma}

\begin{proof}
Let $D_{univ} \subset X \times_S \underline{\Hilbfunctor}^{d + 1}_{X/S}$
be the universal object. There is a commutative diagram
$$
\xymatrix{
\underline{\Hilbfunctor}^d_{X/S} \times_S X \ar[rr] \ar[rd] & &
D_{univ} \ar[ld] \ar@{^{(}->}[r] &
\underline{\Hilbfunctor}^{d + 1}_{X/S} \times_S X \\
& \underline{\Hilbfunctor}^{d + 1}_{X/S}
}
$$
where the top horizontal arrow maps $(D', x)$ to $(D' + x, x)$.
We claim this morphism is an isomorphism
which certainly proves the lemma. Namely, given a scheme $T$ over $S$,
a $T$-valued point $\xi$ of $D_{univ}$ is given by a pair $\xi = (D, x)$
where $D \subset X_T$ is a closed subscheme finite locally free
of degree $d + 1$ over $T$ and $x : T \to X$ is a morphism whose
graph $x : T \to X_T$ factors through $D$. Then by
Lemma \ref{lemma-difference-divisors-on-curves}
we can write $D = D' + x$ for some $D' \subset X_T$ finite locally
free of degree $d$ over $T$. Sending $\xi = (D, x)$ to the pair
$(D', x)$ is the desired inverse.
\end{proof}

\end{document}
